\documentclass[9pt,conference]{IEEEtran}
\usepackage[utf8]{inputenc}
\usepackage[spanish,english]{babel}
\usepackage{graphicx}
\usepackage{booktabs}
\usepackage{caption}
\usepackage{subcaption}
\usepackage{url}
\usepackage{hyperref}
\usepackage{tabularx}
\usepackage{booktabs}
\usepackage{threeparttable}
\usepackage{float}

\title{Food Waste Reduction Platform: Robust System Design and Project Management}

\author{
Luna Alejandra Sandoval Rodríguez \\
\textit{Project Lead \& Systems Analyst} \\
Universidad Distrital Francisco José de Caldas \\
Bogotá, Colombia \\
20241020053
\and
Cristian Camilo Bonilla Lizarazo \\
\textit{Frontend Developer \& Systems Analyst} \\
Universidad Distrital Francisco José de Caldas \\
Bogotá, Colombia \\
20241020015
\and
Nicolás Rodríguez Granados \\
\textit{Backend Developer \& Systems Analyst} \\
Universidad Distrital Francisco José de Caldas \\
Bogotá, Colombia \\
20241020037
\and
Juan Sebastián Bravo Rojas \\
\textit{Database Specialist \& Systems Analyst} \\
Universidad Distrital Francisco José de Caldas \\
Bogotá, Colombia \\
20241020004
}

\begin{document}

\maketitle

\section{Implementation Strategy}
% ============================================================
The Implementation Strategy translates the Workshop 2 technical design and the Workshop 3 experimental validation into a concrete, reproducible development and deployment plan. It follows the best-practice collaboration guidelines established in \textit{TeamWork\_CS.pdf} (feature-branch workflow, atomic commits, mandatory peer reviews, and GitHub Projects board) and maintains full traceability to the 17 design requirements and the four-layer architecture defined in Workshop 2.

\subsection{Implementation Phases}

The Implementation Strategy translates the Workshop 2 technical design and the Workshop 3 experimental validation into a concrete, reproducible development and deployment plan. It follows the best-practice collaboration guidelines established in \textit{TeamWork\_CS.pdf} and maintains full traceability to the 17 design requirements and the four-layer architecture defined in Workshop 2.

\subsection{Implementation Phases}

\begin{enumerate}
    \item \textbf{Phase 1 – Core Implementation and Experimentation (Workshop 3)}  
    Full-stack prototype development using the Workshop 2 technology stack (completed by 23 May 2026).

    \item \textbf{Phase 2 – Simulation, Pilot, and Final MVP (Workshop 4)}  
    Realistic simulations, short controlled pilot, and delivery of all communication artifacts by 30 May 2026.
\end{enumerate}

\subsection{Required Infrastructure and Technology Stack}

The platform is fully containerized to guarantee reproducibility and scalability:

\begin{table}[h]
\centering
\small
\caption{Required Infrastructure and Technology Stack}
\label{tab:infrastructure}
\begin{tabular}{@{}p{1.5cm}p{2.0cm}p{3.5cm}@{}}
\toprule
\textbf{Layer} & \textbf{Component} & \textbf{Technology} \\
\midrule
Frontend & Mobile Application & React Native / Expo (Android \& iOS) \\
Backend & REST APIs \& Business Logic & NestJS (Node.js) \\
Data & Persistent Storage \& Geospatial Queries & PostgreSQL + PostGIS \\
Concurrency & Caching \& Queues & Redis + Message Broker (RabbitMQ/Redis Streams) \\
Notifications & Push Notifications & Firebase Cloud Messaging (FCM) \\
Infrastructure & Containerization \& Orchestration & Docker + Docker Compose \\
Deployment & Cloud Hosting (pilot) & AWS / Google Cloud (free tier for academic use) \\
Monitoring & Analytics \& Dashboards & Custom Analytics Layer (PostgreSQL views) \\
\bottomrule
\end{tabular}
\end{table}

All components are defined in a single \texttt{docker-compose.yml} file stored in the repository root, enabling one-command local deployment for development and testing.

\subsection{Human Resources and Team Allocation}

The implementation is executed by the four-member cross-functional team using the RACI matrix defined in the Project Charter (Section 4):

\begin{itemize}
    \item \textbf{Luna A. Sandoval Rodríguez} (Project Lead) – Overall coordination, integration, and quality gate reviews.
    \item \textbf{Nicolás Rodríguez Granados} (Backend Developer) – Matching algorithm, APIs, and message broker.
    \item \textbf{Cristian C. Bonilla Lizarazo} (Frontend Developer) – Mobile UI/UX and API integration.
    \item \textbf{Juan S. Bravo Rojas} (Database Specialist) – PostGIS schema, data integrity, and analytics layer.
\end{itemize}

Weekly 45–60 minute coordination meetings and mandatory peer code reviews on GitHub ensure knowledge sharing and early detection of integration issues.

\subsection{Estimated Budget and Resources}

As an academic project, the budget is calculated primarily in human-hours with minimal monetary cost:

\begin{table}[h]
\centering
\small
\caption{Estimated Budget (Workshop 3 + Workshop 4)}
\label{tab:budget}
\begin{tabular}{@{}p{2.5cm}p{3.0cm}p{2.5cm}@{}}
\toprule
\textbf{Resource} & \textbf{Description} & \textbf{Cost} \\
\midrule
Human Capital & 4 engineers $\times$ 10 h/week $\times$ 8 remaining weeks & 320 HH (no monetary cost) \\
Cloud Infrastructure & Docker hosting + PostgreSQL + FCM (free tier) & \$0 (academic tier) \\
Production Pilot Hosting & AWS/Google Cloud (estimated for 8-week pilot) & \$50 USD (theoretical) \\
Development Tools & GitHub, Docker, Expo, PostGIS & \$0 (open source) \\
Total Estimated Monetary Cost & & \textbf{\$50 USD} \\
\bottomrule
\end{tabular}
\end{table}

All software components are open-source or free-tier, ensuring zero licensing costs.

\subsection{Rollout Timeline}

The rollout follows a controlled, incremental approach aligned with the course schedule:

\begin{table}[h]
\centering
\small
\caption{Rollout Timeline (Workshop 3 → Final MVP)}
\label{tab:rollout-timeline}
\begin{tabular}{@{}p{1.8cm}p{2.8cm}p{3.5cm}@{}}
\toprule
\textbf{Milestone} & \textbf{Date} & \textbf{Deliverables} \\
\midrule
Core Prototype Complete & 23 May 2026 & Full implementation, unit/integration tests, risk register validation (Workshop 4 delivery) \\
Simulation & 20--23 May 2026 & 300-user load tests, performance benchmarking, sensitivity analysis \\
Pilot Deployment & 24--30 May 2026 & Short controlled validation pilot within 3 km radius (Engineering Faculty area) \\
Final MVP & 30 May 2026 & Production-ready application + administrative dashboard \\
Communication Deliverables & 30 May 2026 & Scientific paper, academic poster, technical report, and final presentation \\
\bottomrule
\end{tabular}
\end{table}

The Gantt chart illustrating the complete timeline is presented in Fig.~\ref{fig:gantt} (Section 4).


\subsection{Resource Management Plan}

\subsubsection{Resource Allocation and Tracking}

Resource usage is tracked through the GitHub Projects Kanban board and weekly burndown charts. Each sprint is limited to a maximum of 10 hours per team member to prevent burnout. Effort distribution is aligned with the RACI matrix defined in the Project Charter:

\begin{table}[h]
\centering
\small
\caption{Resource Allocation Overview (Workshop 3–4)}
\label{tab:resource-allocation}
\begin{tabular}{@{}p{1.8cm}p{2.8cm}p{3.5cm}@{}}
\toprule
\textbf{Resource Type} & \textbf{Management Mechanism} & \textbf{Planned Effort / Cost} \\
\midrule
Human Resources & GitHub Projects + weekly burndown charts & 640 HH total (4 engineers × 10 h/week) \\
Technical Infrastructure & Docker Compose + free-tier cloud services & \$0 (academic tier) \\
Financial Resources & Monitored in project board & \$0--50 (theoretical pilot hosting) \\
\bottomrule
\end{tabular}
\end{table}

\subsubsection{Resource Monitoring and Control}

\begin{itemize}
    \item \textbf{Weekly Resource Review:} During the 45–60 minute status meeting, the team reviews burndown charts and flags any deviation greater than 15\%.
    \item \textbf{Automatic Alerts:} GitHub issue labels (\texttt{resource:critical}, \texttt{resource:overallocated}) trigger notifications to the Project Lead.
    \item \textbf{Control Gates:} No merge to \texttt{develop} is allowed without peer review and confirmation that the assigned effort was not exceeded.
    \item \textbf{Contingency:} If any member exceeds 12 hours in a week, the Project Lead reassigns tasks immediately to maintain balance.
\end{itemize}

This plan guarantees that the project remains within academic constraints, ensures full traceability of effort, and supports the successful delivery of the final MVP and communication artifacts by \textbf{30 May 2026}.



\begin{thebibliography}{10}
\bibitem{w1} Workshop 1 – Food Waste Reduction Mobile Application, A Rigorous System Analysis, Universidad Distrital Francisco José de Caldas, 2026.
\bibitem{w2} Workshop 2 – System Design Document, Universidad Distrital Francisco José de Caldas, 2026.
\bibitem{ieee} IEEE Standards Association, ``IEEE 730-2014 - IEEE Standard for Software Quality Assurance Processes,'' 2014.
\end{thebibliography}

\end{document}
